\documentclass[10pt,a4paper]{amsart}
\usepackage{amsmath,amssymb,amsthm,graphicx}
% For commutative diagrams we use Xy-pic
\usepackage[all]{xy}

% We use 2cell for 2-commutative diagrams.
\xyoption{2cell}
% Theorem environments.
%
\theoremstyle{plain}
\newtheorem{theorem}[subsection]{Theorem}
\newtheorem{proposition}[subsection]{Proposition}
\newtheorem{lemma}[subsection]{Lemma}

\theoremstyle{definition}
\newtheorem{definition}[subsection]{Definition}
\newtheorem{example}[subsection]{Example}
\newtheorem{exercise}[subsection]{Exercise}
\newtheorem{situation}[subsection]{Situation}

\theoremstyle{remark}
\newtheorem{remark}[subsection]{Remark}
\newtheorem{remarks}[subsection]{Remarks}

\numberwithin{equation}{subsection}

% Macros
%
\def\lim{\mathop{\mathrm{lim}}\nolimits}
\def\colim{\mathop{\mathrm{colim}}\nolimits}
\def\Spec{\mathop{\mathrm{Spec}}}
\def\Hom{\mathop{\mathrm{Hom}}\nolimits}
\def\Ext{\mathop{\mathrm{Ext}}\nolimits}
\def\SheafHom{\mathop{\mathcal{H}\!\mathit{om}}\nolimits}
\def\SheafExt{\mathop{\mathcal{E}\!\mathit{xt}}\nolimits}
\def\Sch{\mathit{Sch}}
\def\Mor{\mathop{\mathrm{Mor}}\nolimits}
\def\Ob{\mathop{\mathrm{Ob}}\nolimits}
\def\Sh{\mathop{\mathit{Sh}}\nolimits}
\def\NL{\mathop{N\!L}\nolimits}
\def\CH{\mathop{\mathrm{CH}}\nolimits}
\def\proetale{{pro\text{-}\acute{e}tale}}
\def\etale{{\acute{e}tale}}
\def\QCoh{\mathit{QCoh}}
\def\Ker{\mathop{\mathrm{Ker}}}
\def\Im{\mathop{\mathrm{Im}}}
\def\Coker{\mathop{\mathrm{Coker}}}
\def\Coim{\mathop{\mathrm{Coim}}}

% Boxtimes
%
\DeclareMathSymbol{\boxtimes}{\mathbin}{AMSa}{"02}

%
% Macros for moduli stacks/spaces
%
\def\QCohstack{\mathcal{QC}\!\mathit{oh}}
\def\Cohstack{\mathcal{C}\!\mathit{oh}}
\def\Spacesstack{\mathcal{S}\!\mathit{paces}}
\def\Quotfunctor{\mathrm{Quot}}
\def\Hilbfunctor{\mathrm{Hilb}}
\def\Curvesstack{\mathcal{C}\!\mathit{urves}}
\def\Polarizedstack{\mathcal{P}\!\mathit{olarized}}
\def\Complexesstack{\mathcal{C}\!\mathit{omplexes}}
% \Pic is the operator that assigns to X its picard group, usage \Pic(X)
% \Picardstack_{X/B} denotes the Picard stack of X over B
% \Picardfunctor_{X/B} denotes the Picard functor of X over B
\def\Pic{\mathop{\mathrm{Pic}}\nolimits}
\def\Picardstack{\mathcal{P}\!\mathit{ic}}
\def\Picardfunctor{\mathrm{Pic}}
\def\Deformationcategory{\mathcal{D}\!\mathit{ef}}
\usepackage{cleveref}
\crefname{originallemma}{original lemma}{original lemmas}
\Crefname{originallemma}{Original lemma}{Original lemmas}
\makeatletter
\newlabel{lemma-dimension-at-most-one}{{02P7}{source}{}{sourcepointer-1}{}}
\newlabel{lemma-dimension-at-most-one@cref}{{[originallemma][0][]02P7}{[1][0][]source}{}{}{}}
\makeatother
\crefname{originallemma}{original lemma}{original lemmas}
\Crefname{originallemma}{Original lemma}{Original lemmas}
\makeatletter
\newlabel{lemma-determinant-dimension-one}{{02P9}{source}{}{sourcepointer-2}{}}
\newlabel{lemma-determinant-dimension-one@cref}{{[originallemma][0][]02P9}{[1][0][]source}{}{}{}}
\makeatother
\title{A canonical isomorphism between the determinant of the outer terms and the determinant of the middle term of a short exact sequence of finite length modules, with three compatibility properties}
\author{The Stacks Project Authors}
\date{stacks-project:chow.tex:local-snapshot-20261012; GNU FDL 1.2 or later}
\begin{document}
\maketitle
\noindent\textit{Editorial note.} All delivered mathematical text is the author's own, copied from the pinned Stacks Project chapter file chow.tex; no statement or proof step was rewritten, reordered or removed.\par\medskip
\noindent\textit{Editorial note.} Preamble provenance: the pinned chapter file opens with an input of the shared preamble and carries no preamble text of its own, so the project's single Stacks-wide preamble staged at sources/stacks-chow/preamble.tex is excerpted verbatim as two spans, lines 14-18 (Xy-pic setup) and lines 154-217 (the native theorem declarations and mathematical macros). Nothing was added, removed or redefined; the wrapper class and the amsmath/amssymb/amsthm/graphicx base packages are supplied by the assembler.\par\medskip
\noindent\textit{Editorial note.} The counted statement and proof contain no cross reference, no citation and no included file, so the delivered body is self contained.\par\medskip
\noindent\textit{Editorial note.} Canonical location of the counted proposition: Stacks tag 02PA (sources/stacks-chow/pinned-original-source.tex line 16203).\par\medskip
\noindent\textit{Allowed standard prerequisites (editorial).}\begin{itemize}
\item Admissible sequences, their relations, and the generation of det\_kappa by a single symbol. Same chapter, chow.tex, Section on periodic complexes and determinants, in particular the relation (c) used in the sign computation at proof lines 16354-16362. Supplies the three admissible relations of det\_kappa and the presentation of det\_kappa by symbols; the transfer of relations and the ordered application of relation (c) are carried out in the counted span.
\item Generation of det\_kappa of a finite length module by one appropriate symbol. Same chapter, Lemma labelled lemma-determinant-dimension-one (source line 15850) and Lemma labelled lemma-dimension-at-most-one (source line 15761), both invoked at proof lines 16306 and 16322. Supplies that det\_kappa(L) is generated by any single symbol attached to an admissible sequence generating L, which is what makes gamma surjective; the construction of gamma and the verification of its three properties stay inside the counted span.
\end{itemize}\medskip
\ifdefined\hypertarget\hypertarget{sourcepointer-1}{}\fi
\noindent\textit{Source reference.} Original lemma 02P7: sources/stacks-chow/pinned-original-source.tex, source lines 15761--15761. Same-chapter pointer into the author environment lemma at source line 15761; invoked at proof line 16307 of the counted proof. It is retained only as a navigation pointer; every load-bearing step is inside the delivered span.\par\medskip
\ifdefined\hypertarget\hypertarget{sourcepointer-2}{}\fi
\noindent\textit{Source reference.} Original lemma 02P9: sources/stacks-chow/pinned-original-source.tex, source lines 15850--15850. Same-chapter pointer into the author environment lemma at source line 15850; invoked at proof line 16322 of the counted proof. It is retained only as a navigation pointer; every load-bearing step is inside the delivered span.\par\medskip
\begin{lemma}
\label{lemma-det-exact-sequences}
Let $(R, \mathfrak m, \kappa)$ be a local ring.
For every short exact sequence
$$
0 \to K \to L \to M \to 0
$$
of finite length $R$-modules there exists a canonical isomorphism
$$
\gamma_{K \to L \to M} :
\det\nolimits_\kappa(K) \otimes_\kappa \det\nolimits_\kappa(M)
\longrightarrow
\det\nolimits_\kappa(L)
$$
defined by the rule on nonzero symbols
$$
[e_1, \ldots, e_k]
\otimes
[\overline{f}_1, \ldots, \overline{f}_m]
\longrightarrow
[e_1, \ldots, e_k, f_1, \ldots, f_m]
$$
with the following properties:
\begin{enumerate}
\item For every isomorphism of short exact sequences, i.e., for
every commutative diagram
$$
\xymatrix{
0 \ar[r] &
K \ar[r] \ar[d]^u &
L \ar[r] \ar[d]^v &
M \ar[r] \ar[d]^w &
0 \\
0 \ar[r] &
K' \ar[r] &
L' \ar[r] &
M' \ar[r] &
0
}
$$
with short exact rows and isomorphisms $u, v, w$ we have
$$
\gamma_{K' \to L' \to M'} \circ
(\det\nolimits_\kappa(u) \otimes \det\nolimits_\kappa(w))
=
\det\nolimits_\kappa(v) \circ
\gamma_{K \to L \to M},
$$
\item for every commutative square of finite length $R$-modules
with exact rows and columns
$$
\xymatrix{
& 0 \ar[d] & 0 \ar[d] & 0 \ar[d] & \\
0 \ar[r] & A \ar[r] \ar[d] & B \ar[r] \ar[d] & C \ar[r] \ar[d] & 0 \\
0 \ar[r] & D \ar[r] \ar[d] & E \ar[r] \ar[d] & F \ar[r] \ar[d] & 0 \\
0 \ar[r] & G \ar[r] \ar[d] & H \ar[r] \ar[d] & I \ar[r] \ar[d] & 0 \\
& 0  & 0  & 0  &
}
$$
the following diagram is commutative
$$
\xymatrix{
\det\nolimits_\kappa(A) \otimes
\det\nolimits_\kappa(C) \otimes
\det\nolimits_\kappa(G) \otimes
\det\nolimits_\kappa(I)
\ar[dd]_{\epsilon}
\ar[rrr]_-{\gamma_{A \to B \to C} \otimes \gamma_{G \to H \to I}}
& & &
\det\nolimits_\kappa(B) \otimes
\det\nolimits_\kappa(H)
\ar[d]^{\gamma_{B \to E \to H}}
\\
& & & \det\nolimits_\kappa(E)
\\
\det\nolimits_\kappa(A) \otimes
\det\nolimits_\kappa(G) \otimes
\det\nolimits_\kappa(C) \otimes
\det\nolimits_\kappa(I)
\ar[rrr]^-{\gamma_{A \to D \to G} \otimes \gamma_{C \to F \to I}}
& & &
\det\nolimits_\kappa(D) \otimes
\det\nolimits_\kappa(F)
\ar[u]_{\gamma_{D \to E \to F}}
}
$$
where $\epsilon$ is the switch of the factors in the tensor product
times $(-1)^{cg}$ with $c = \text{length}_R(C)$ and $g = \text{length}_R(G)$,
and
\item the map $\gamma_{K \to L \to M}$ agrees with the usual isomorphism
if $0 \to K \to L \to M \to 0$ is actually a short exact sequence
of $\kappa$-vector spaces.
\end{enumerate}
\end{lemma}

\begin{proof}
The significance of taking nonzero symbols in the explicit description
of the map $\gamma_{K \to L \to M}$ is simply that if $(e_1, \ldots, e_l)$
is an admissible sequence in $K$, and
$(\overline{f}_1, \ldots, \overline{f}_m)$ is an admissible sequence in
$M$, then it is not guaranteed that $(e_1, \ldots, e_l, f_1, \ldots, f_m)$
is an admissible sequence in $L$ (where of course $f_i \in L$ signifies
a lift of $\overline{f}_i$). However, if the symbol
$[e_1, \ldots, e_l]$ is nonzero in $\det_\kappa(K)$, then
necessarily $K = \langle e_1, \ldots, e_k\rangle$ (see
proof of Lemma \ref{lemma-dimension-at-most-one}), and
in this case it is true that $(e_1, \ldots, e_k, f_1, \ldots, f_m)$
is an admissible sequence.
Moreover, by the admissible relations of type (b) for $\det_\kappa(L)$
we see that the value of $[e_1, \ldots, e_k, f_1, \ldots, f_m]$ in
$\det_\kappa(L)$ is independent of the choice of the lifts
$f_i$ in this case also. Given this remark, it is clear
that an admissible relation for $e_1, \ldots, e_k$ in $K$
translates into an admissible relation among
$e_1, \ldots, e_k, f_1, \ldots, f_m$ in $L$, and
similarly for an admissible relation among the
$\overline{f}_1, \ldots, \overline{f}_m$.
Thus $\gamma$ defines a linear map of vector spaces as claimed in the lemma.

\medskip\noindent
By Lemma \ref{lemma-determinant-dimension-one} we know
$\det_\kappa(L)$ is generated by any single
symbol $[x_1, \ldots, x_{k + m}]$ such that
$(x_1, \ldots, x_{k + m})$ is an admissible sequence
with $L = \langle x_1, \ldots, x_{k + m}\rangle$. Hence it is
clear that the map $\gamma_{K \to L \to M}$ is surjective and
hence an isomorphism.

\medskip\noindent
Property (1) holds because
\begin{eqnarray*}
& & \det\nolimits_\kappa(v)([e_1, \ldots, e_k, f_1, \ldots, f_m]) \\
& = &
[v(e_1), \ldots, v(e_k), v(f_1), \ldots, v(f_m)] \\
& = &
\gamma_{K' \to L' \to M'}([u(e_1), \ldots, u(e_k)]
\otimes [w(f_1), \ldots, w(f_m)]).
\end{eqnarray*}
Property (2) means that given a symbol
$[\alpha_1, \ldots, \alpha_a]$ generating $\det_\kappa(A)$,
a symbol $[\gamma_1, \ldots, \gamma_c]$ generating $\det_\kappa(C)$,
a symbol $[\zeta_1, \ldots, \zeta_g]$ generating $\det_\kappa(G)$, and
a symbol $[\iota_1, \ldots, \iota_i]$ generating $\det_\kappa(I)$
we have
\begin{eqnarray*}
& & [\alpha_1, \ldots, \alpha_a, \tilde\gamma_1, \ldots, \tilde\gamma_c,
\tilde\zeta_1, \ldots, \tilde\zeta_g, \tilde\iota_1, \ldots, \tilde\iota_i] \\
& = &
(-1)^{cg} [\alpha_1, \ldots, \alpha_a, \tilde\zeta_1, \ldots, \tilde\zeta_g,
\tilde\gamma_1, \ldots, \tilde\gamma_c, \tilde\iota_1, \ldots, \tilde\iota_i]
\end{eqnarray*}
(for suitable lifts $\tilde{x}$ in $E$) in $\det_\kappa(E)$.
This holds because we may use the admissible relations of type (c)
$cg$ times in the following order: move the
$\tilde\zeta_1$ past the elements
$\tilde\gamma_c, \ldots, \tilde\gamma_1$
(allowed since $\mathfrak m\tilde\zeta_1 \subset A$),
then move $\tilde\zeta_2$ past the elements
$\tilde\gamma_c, \ldots, \tilde\gamma_1$
(allowed since $\mathfrak m\tilde\zeta_2 \subset A + R\tilde\zeta_1$),
and so on.

\medskip\noindent
Part (3) of the lemma is obvious.
This finishes the proof.
\end{proof}

\end{document}
